\section{Identification of Embodiment and AR Effects}
\subsection{Causal contrasts}
The three-condition design identifies the incremental contrast of movement without AR (E--S) and AR added to movement (A--E), subject to successful task matching. It is not a full factorial decomposition of every possible movement-by-AR interaction. The total A--S contrast alone cannot isolate AR. All conditions will share the same cognitive rule, challenge exposure, reward schedule, and recovery definition to the extent technically feasible.

\subsection{Task equivalence gate}
Before mechanism testing, pilot work will compare comprehension, baseline error rates, effective response opportunities, input latency, and completion burden. A mismatch that changes the cognitive construct requires task redesign rather than statistical adjustment alone. Movement quantity, spatial tracking quality, and device capability are reported as manipulation and fidelity variables, not automatically interpreted as improvements.

\subsection{Non-inferiority and transfer}
Non-inferiority margins for cognitive accuracy and other safety-critical performance measures require independent justification before outcome access. Both intention-to-assign and technically adherent summaries should be examined because missing AR tracking may bias an apparently favorable complete-case estimate. Independent transfer will be tested in a subsequent fixed-context versus varied-context randomized study with matched dose. Its primary endpoint must be measured outside the trained AR interface, with the follow-up interval and missing-data strategy frozen prospectively.

\subsection{Falsification and feasibility}
Greater physical movement without preserved executive-task performance does not support a beneficial embodiment mechanism. A--S improvement without A--E improvement cannot be attributed specifically to AR. Effects restricted to novel first exposure, high-end devices, or well-tracked events do not justify a general BYOD claim. Safety interruptions, visual discomfort, fatigue, and environmental limitations remain outcomes even when scientific recovery data are invalid.